\documentclass[oneside,a4paper,14pt]{extarticle}
\usepackage[a4paper,letterpaper,top=20mm,bottom=20mm,left=20mm,right=10mm]{geometry}
\usepackage[russian]{babel}
\usepackage{indentfirst}
\usepackage{graphicx}
\usepackage{caption}
\usepackage{titlesec}
\usepackage{minted, fancyvrb}
\usepackage{hyperref}
\usepackage{enumitem}
\usepackage{amsmath}
\usepackage{fancyhdr}
\usepackage{colortbl}
\usepackage[dvipsnames]{xcolor}
\usepackage{float}

% Форматирование листингов с кодом
\setminted{style = rainbow_dash, fontsize = \small}

% Форматирование заголовков
\titleformat{\section}{\normalsize\bfseries}{\thesection}{1em}{}
\titleformat{\subsection}{\normalsize\bfseries}{\thesubsection}{1em}{}
\titleformat{\subsubsection}{\normalsize\bfseries}{\thesubsubsection}{1em}{}

% Интерлиньяж и абзац
\renewcommand\baselinestretch{1.33}
\setlength{\parindent}{1.25cm}

% Для всех списков
\setlist[enumerate]{
  left=0.5\parindent,
  label=\arabic*.,
  itemsep=0pt,
  topsep=5pt,
  partopsep=0pt,
  parsep=0pt
}
\setlist[itemize]{
  left=0.5\parindent,
  itemsep=0pt,
  topsep=5pt,
  partopsep=0pt,
  parsep=0pt
}

% Гиперссылки
\hypersetup{
  colorlinks=true,
  linkcolor=black,
  urlcolor=blue,
  pdfborder={0 0 0},
  pdftitle={Построение 2-разрядного регистра на D-триггерах (ЛА3, ЛА4)},
  pdfauthor={Черкасов А.А. и Зинин В.А.},
}

\begin{document}

\newpage
\thispagestyle{empty}
\begin{center}
  МИНИСТЕРСТВО НАУКИ И ВЫСШЕГО ОБРАЗОВАНИЯ РОССИЙСКОЙ ФЕДЕРАЦИИ\\
  ФЕДЕРАЛЬНОЕ ГОСУДАРСТВЕННОЕ БЮДЖЕТНОЕ УЧРЕЖДЕНИЕ ВЫСШЕГО ОБРАЗОВАНИЯ\\
  <<ВЯТСКИЙ ГОСУДАРСТВЕННЫЙ УНИВЕРСИТЕТ>>\\
  Институт математики и информационных систем\\
  Факультет автоматики и вычислительной техники\\
  Кафедра электронных вычислительных машин
\end{center}
\vspace{10mm}

\hfill
\begin{tabular}{l}
  \footnotesize Дата сдачи на проверку:                                          \\
  \footnotesize <<\rule[-1mm]{5mm}{0.10mm}\/>>\rule[-1mm]{20mm}{0.10mm}\ 2025 г. \\
  \footnotesize Проверено:                                                       \\
  \footnotesize <<\rule[-1mm]{5mm}{0.10mm}\/>>\rule[-1mm]{20mm}{0.10mm}\ 2025 г. \\
\end{tabular}
\vfill

\begin{center}
  Отчёт по лабораторной работе №2\\
  по дисциплине\\
  <<Цифровая схемотехника>>\\
  \textbf{Построение 2-разрядного регистра на D-триггерах}\\
\end{center}
\vspace{25mm}
\noindent
\begin{tabular}{ll}
  Разработали студенты гр. ИВТб-2301-05-00 & \hspace{5mm}\rule[-1mm]{30mm}{0.10mm}\,/Черкасов А. А./ \\
                                           & \hspace{5mm}\rule[-1mm]{30mm}{0.10mm}\,/Зинин В. А./    \\
                                           & \hspace{12.5mm}\footnotesize(подпись)                   \\
  Преподаватель                           & \hspace{5mm}\rule[-1mm]{30mm}{0.10mm}\,/Мельцов В. Ю./  \\
                                           & \hspace{12.5mm}\footnotesize(подпись)                   \\
\end{tabular}

\noindent
\begin{tabular}{lp{58mm}r}
  Работа защищена &  & \hspace{5mm}<<\rule[-1mm]{5mm}{0.10mm}\/>>\rule[-1mm]{30mm}{0.10mm}\ 2026 г.
\end{tabular}
\vfill

\begin{center}
  Киров\\
  2026
\end{center}

\newpage\thispagestyle{plain}

\section{Цель лабораторной работы}
Получить навыки построения и анализа 2-разрядных регистров на D-триггерах. Изучить принцип параллельной записи данных. Выполнить синтез функциональной и принципиальной схем регистра на базе логических элементов ЛА3 и ЛА4 (серия ТТЛ).

\section{Теоретическое описание}
\subsection{D-триггер}
D-триггер (триггер задержки) имеет информационный вход $D$, тактовый вход $C$ и асинхронные входы установки $\overline{S}$ и сброса $\overline{R}$ (активный низкий уровень). Переключение состояния триггера по тактовому входу происходит по спаду импульса (символ $↓$ в таблице). В данной работе D-триггеры построены на основе элементов И-НЕ (микросхемы ЛА3, ЛА4).

Таблица истинности D-триггера с асинхронными входами:

\begin{center}
  Таблица 1 --- Таблица истинности D-триггера \\
  \vspace{5mm}
  \begin{tabular}{|c|c|c|c|c|c|}
    \hline
    $\overline{S}$ & $\overline{R}$ & C    & D    & $Q_{n+1}$ & $\overline{Q}_{n+1}$ \\
    \hline
    0    & 1    & X    & X    & 1    & 0    \\
    1    & 0    & X    & X    & 0    & 1    \\
    0    & 0    & X    & X    & 1*   & 1*   \\
    1    & 1    & $\downarrow$ & 1 & 1    & 0    \\
    1    & 1    & $\downarrow$ & 0 & 0    & 1    \\
    1    & 1    & 0    & X    & $Q_n$ & $\overline{Q}_n$ \\
    \hline
  \end{tabular}
  \vspace{2mm}
  \par\footnotesize * --- неустойчивое состояние.
\end{center}

\subsection{2-разрядный параллельный регистр}
Регистр предназначен для хранения двух бит информации. Каждый разряд представляет собой отдельный D-триггер. Информационный вход $D_i$ подаётся непосредственно на $D$-вход соответствующего триггера. Тактовые входы всех триггеров объединяются и подключаются к общему сигналу синхронизации $C$. По спаду тактового импульса происходит запись данных.

\section{Функциональная схема регистра}
На рисунке 1 представлена функциональная схема 2-разрядного параллельного регистра. Вход $D$ является информационным, выход $Q_i$ — выходом регистра. Все триггеры тактируются общим сигналом $C$.

\begin{figure}[H]
  \centering
  \includegraphics[width=0.9\textwidth]{pics/functional_diagram.png}
  \caption*{Рисунок 1 --- Функциональная схема 2-разрядного регистра на D-триггерах}
\end{figure}

\section{Принципиальная схема устройства}
Принципиальная схема (рисунок 2) выполнена на микросхемах серии ТТЛ: \textbf{ЛА3} и \textbf{ЛА4}. ЛА3 (К155ЛА3) содержит четыре элемента 2И-НЕ, ЛА4 (К155ЛА4) — три элемента 3И-НЕ. Из этих элементов собраны два D-триггера, образующих 2-разрядный регистр. Схема включения микросхем соответствует заводским цоколёвкам. Питание подаётся от источника +5~В.

\begin{figure}[H]
  \centering
  \includegraphics[width=0.9\textwidth]{pics/principle_diagram.png}
  \caption*{Рисунок 2 --- Принципиальная схема 2-разрядного регистра на D-триггерах (ЛА3, ЛА4)}
\end{figure}

\newpage
\section{Таблица питания микросхем}
В соответствии с технической документацией на микросхемы серии ТТЛ (ЛА3, ЛА4) напряжение питания составляет +5~В. Выводы питания для всех использованных микросхем одинаковы.

\begin{center}
  Таблица 2 --- Таблица питания для микросхем ТТЛ \\
  \vspace{5mm}
  \begin{tabular}{|c|c|c|}
    \hline
    Микросхемы     & +5V & 0V \\
    \hline
    D1, D2 & 14 & 07   \\
    \hline
  \end{tabular}
\end{center}
\vspace{5mm}

\noindent
Здесь D1 и D2 обозначают микросхемы ЛА4 и ЛА3 соответственно.

\section*{Вывод}
В ходе лабораторной работы были получены практические навыки построения 2-разрядного параллельного регистра на D-триггерах. Триггеры реализованы на базе элементов И-НЕ (микросхемы ЛА3 и ЛА4). Разработаны функциональная и принципиальная схемы устройства, указаны цепи питания. Полученные результаты могут быть использованы при проектировании более сложных регистров, счётчиков и блоков памяти на ТТЛ-логике.

\end{document}